\chapter{爵士：把对话写进和声里}
\label{ch:genre-jazz-deep}

\classquote{大乐必易，大礼必简。}{《礼记·乐记》}{真正宏大的音乐以平易为底色；隆重的仪式以简约为本。爵士的表面繁复，底下是极简的对话规则。}

\section{什么是爵士：先听它“怎么一起响”}
\label{sec:jazz-what}

\term{爵士}{jazz}不是一种音阶，也不是“复杂的和弦”。它更像一套关于“怎样一起演奏”的约定：一个人给出主题，另一个人接着回答；节奏组在底下稳稳地走，前台在和弦上现场造句。录音里那些最动人的瞬间，常常不是写好的音符，而是几个人在同一拍上做出的同一个决定。

把爵士听懂，先抓四件事：

\begin{longtable}{p{0.16\textwidth}p{0.32\textwidth}p{0.38\textwidth}}
\toprule
抓手 & 听感 & 写作问题 \\
\midrule
摇摆 & 八分音符“长—短”不均等，身体想跟着晃 & 这处 Swing 是均匀的，还是真的把后半拍往后推了？ \\
行走贝斯 & 每拍一个音，像在地面铺台阶 & 贝斯只有根音，还是有经过音把和弦连起来？ \\
问答 & 一句先问、一句后答，像两个人说话 & Comping 与独奏有没有留出“回答”的空？ \\
即兴 & 主题之后展开新句子，每次都不同 & Solo 是在和弦上发展动机，只是把音阶跑了一遍？ \\
\bottomrule
\end{longtable}

爵士最怕的不是“不标准”，而是“没有对话”。一段只有和弦进行、没有呼应的独奏，像一个人对着墙说话。爵士的专业感来自互相倾听——节奏组听见独奏者往前推，就跟着推；独奏者听见留白，就往空白处填。

\section{从新奥尔良到全球：爵士的来路}
\label{sec:jazz-history}

爵士一路吸收布鲁斯、拉格泰姆、进行曲、福音与拉丁节奏，再把它们整理成“可即兴”的语言。下表只抓“什么场合改变了演奏方式，配器就跟着变”的主线。

\begin{longtable}{p{0.16\textwidth}p{0.34\textwidth}p{0.42\textwidth}}
\toprule
时期 & 关键变化 & 对编配的影响 \\
\midrule
\endfirsthead
\toprule
时期 & 关键变化 & 对编配的影响 \\
\midrule
\endhead
1890—1910 年代 & 新奥尔良的铜管乐队、葬礼与游行；布鲁斯与拉格泰姆 & 小号/长号/单簧管前台即兴，节奏组托底；集体即兴多于独奏 \\
1920 年代 & 爵士时代、唱片与舞厅；芝加哥与纽约 & 小乐队编制定型：小号、单簧管/萨克斯、钢琴、贝斯、鼓 \\
1930—1940 年代 & 大乐队（Big Band）与摇摆乐 & 编制扩为若干小号/长号/萨克斯加节奏组，分声部写作；Arranger 成为核心角色 \\
1940 年代 & 比波普（Bebop） & 速度更快、和弦更密；编制缩回小乐队，突出独奏者的技巧 \\
1950 年代 & 酷爵士、硬波普与调式爵士 & 音色更收敛或更蓝调化；\term{调式}{modal}思路减少和弦密度，给独奏更多空间 \\
1960 年代 & 自由爵士与融合爵士的萌芽 & 结构更松，电声乐器进入；节拍与调性边界被重新讨论 \\
1970—1990 年代 & 融合（Fusion）、电声与跨界 & 电贝斯、合成器与摇滚节奏进入；爵士和声与流行编制混用 \\
2000 年代至今 & 全球巡演、流媒体与跨 genres 合作 & 年轻音乐家同时熟悉爵士、嘻哈、电子与古典；编制与曲风边界更松 \\
\bottomrule
\end{longtable}

爵士也是一门“迁移”的语言：它到了不同地方就和当地音乐结合——巴西的 Bossa Nova、古巴的 Afro-Cuban、欧洲的即兴剧场、日本的爵士咖啡馆文化。对中国听众而言，爵士常先从“咖啡馆里的钢琴三重奏”被认识，再慢慢听到它的布鲁斯根与即兴内核。

\section{配器：节奏组与前台的分工}
\label{sec:jazz-craft}

爵士的编制可以用一句话记住：\textbf{节奏组负责时间与和声，前台负责旋律与故事}。把这句话写进轨道名，比堆十条音色更有效。

\subsection{典型编制与两层角色}
\label{sec:jazz-instrumentation}

\begin{longtable}{p{0.16\textwidth}p{0.25\textwidth}p{0.25\textwidth}p{0.20\textwidth}}
\toprule
功能层 & 常见乐器/轨道 & 主要任务 & 常见错误 \\
\midrule
\endfirsthead
\toprule
功能层 & 常见乐器/轨道 & 主要任务 & 常见错误 \\
\midrule
\endhead
节奏组 & 钢琴/吉他、贝斯、鼓 & 定时间、走和声、留空间 & 三件乐器都去抢同一个音区 \\
前台 & 小号、萨克斯、长号、人声 & 给主题、做独奏、讲故事 & 独奏从头响到尾，不给别人留白 \\
色彩 & Vibraphone、弦乐、管风琴 & 增厚织体、改变段落颜色 & 色彩层盖住 Walking Bass 的低频 \\
\bottomrule
\end{longtable}

最简的起点是四件：钢琴、贝斯、鼓，加一件前台旋律（小号或萨克斯）。先把“节奏组＋前台”的关系跑通，再决定要不要加第二管或 Vibraphone。轨数不是爵士深度的证明；会不会对话才是。

\subsection{节奏：Swing 不是“写长写短”}
\label{sec:jazz-groove}

\term{摇摆}{swing}是把八分音符的后半拍往后推，形成“长—短”的律动。关键认知：Swing 是\textbf{演奏感}，不是把音符时值机械地写成三连音。在夜半里可以用量化摆动（swing quantize）或手动把后半拍往后挪一点来做，但更重要的听感来自 Ride 镲与鼓手的轻响分布：

\begin{itemize}
\item \textbf{Ride 镲}：持续走“叮—叮—哒”的摇摆型，是爵士时间的钟。
\item \textbf{军鼓}：轻点 2、4 拍，不追求摇滚式的重击。
\item \textbf{底鼓}：稀疏呼应（feathering），给旋律让路。
\item \textbf{Comping}：钢琴/吉他在关键拍给和弦，其余留空——像“问答”里的答句。
\end{itemize}

如果整段每一拍都同样满，摇摆就没有弹性。节奏组的专业不是“每拍都打”，而是“该走的时候稳得住，该留的时候真的空”。

\subsection{和声：ii–V–I 与调式的两种空间}
\label{sec:jazz-harmony}

爵士和声的核心语法是 \term{ii–V–I}{ii-V-I}：二级小七和弦、五级属七和弦、一级大七和弦的解决。布鲁斯则用 I–IV–V，并允许\term{蓝调音}{blue-note}（降三、降五、降七的弯曲音）。两种语法给出不同的写作空间：

\begin{itemize}
\item \textbf{功能和声（bebop 等）}：和弦密度高，独奏要“追着变化跑”，适合练速度与控制。
\item \textbf{调式爵士（modal）}：和弦少、一个调式停留久，独奏可以“在一个音区里做文章”，适合练动机发展与音色。
\end{itemize}

对夜半里的写作，一个实用建议：先用调式爵士练“少和弦、长线条”，再用 ii–V–I 练“快和弦、短动机”。前者教你发展主题，后者教你跟上变化。

\subsection{旋律：主题要短，独奏要长}
\label{sec:jazz-melody}

爵士的\term{主题}{head}通常是 8、16 或 32 小节的可记短句；真正展开的是之后的\term{即兴}{improvisation}。写主题时记住两件事：

\begin{enumerate}
\item 主题要短到能被记住、能被变奏——别写成流行歌的长句。
\item 主题的句子最好留一个“可发展”的动机（一个节奏型或两三个音），让独奏能拿它做文章。
\end{enumerate}

独奏不是把音阶从上往下跑一遍。好的独奏会重复一个动机、把它移到新的和弦上、改变它的节奏，再在结尾回到主题——听众才能在“新句子”里听见“旧主题”的影子。

\subsection{声音设计：让独奏有地方站}
\label{sec:jazz-sound}

爵士的“编配”很大一部分是\textbf{动态}：谁在响、谁退后。三条可执行原则：

\begin{itemize}
\item \textbf{独奏进入时节奏组降密度}：钢琴少给和弦，鼓把重击换成 Ride 与轻军鼓，贝斯保持走动但不加花。
\item \textbf{问答要真的留空}：Comping 落一个和弦之后，给独奏者半拍到一拍的回应空间。
\item \textbf{低频只留一个主角}：行走贝斯占低频，钢琴左手与贝斯同步时不要另加低音；Vibraphone 的低音杆也要让位。
\end{itemize}

\section{四首曲子的配器拆解}
\label{sec:jazz-cases}

下面四首只做编配分析，不引歌词、不附旋律谱例。作品信息以专辑发行资料与官方页面为据；链接见章末扩展阅读。

\subsection{《So What》：少即是多的调式范本}
\label{sec:jazz-case-so-what}

Miles Davis 的《So What》收录于 1959 年专辑《Kind of Blue》。它是调式爵士的代表：全曲几乎只在两个调式之间来回，给独奏者留下极大的空间。

\begin{itemize}
\item \textbf{和声}：不是密集的 ii–V–I，而是长时间停留在一个调式上；独奏者可以“横向”发展动机，而不是“纵向”追和弦。
\item \textbf{主题}：一个极简的下行动机，由小号先陈述，萨克斯回应——主题本身就是一段问答。
\item \textbf{节奏组}：钢琴 Comping 稀疏，贝斯行走平稳，鼓用 Ride 维持时间；整首歌“退”得很干净。
\item \textbf{编配逻辑}：用减法制造空间，让每个独奏者的句子都听得清楚。
\end{itemize}

可借鉴：如果你的 Solo 总是像跑音阶，试试把和弦减到两个、停留时间拉长，逼自己用动机发展而不是和弦切换。

\subsection{《Take the 'A' Train》：大乐队的声部写作}
\label{sec:jazz-case-a-train}

Duke Ellington 与 Billy Strayhorn 的《Take the 'A' Train》是 1940 年代大乐队摇摆乐的标志曲。它示范了“编曲”怎样把一组乐器写成一个有层次的整体。

\begin{itemize}
\item \textbf{编制}：小号、长号、萨克斯分成声部，与节奏组对话；不同声部轮流承担主题与应答。
\item \textbf{音色}：用弱音器与音区区分声部，让“谁在说话”一目了然。
\item \textbf{节奏}：摇摆的推进来自整个乐队的重音对齐，而不是某一个人加花。
\item \textbf{编配逻辑}：先写声部功能，再写旋律；每一组都有明确任务。
\end{itemize}

可借鉴：即使你只有五条 MIDI 轨，也可以模仿大乐队的“声部感”——把乐器分组，让组与组之间问答，而不是所有人同时演奏同一条旋律。

\subsection{《My Favorite Things》：把主题拆成独奏的积木}
\label{sec:jazz-case-favorite-things}

John Coltrane 1960 年录制的《My Favorite Things》展示了怎样把一条简单主题变成长篇独奏的素材。

\begin{itemize}
\item \textbf{主题}：一条易记的旋律，结构规整；它既是主题，也是独奏的“积木”。
\item \textbf{独奏}：Coltrane 把主题动机拆开、移调、加密，再重新组装；听众始终能听见主题的影子。
\item \textbf{节奏组}：钢琴与贝斯在长独奏中保持稳定循环，给独奏者一个可依赖的地面。
\item \textbf{编配逻辑}：不是“主题 + 随便即兴”，而是“主题 → 主题的变形 → 主题回归”。
\end{itemize}

可借鉴：写 Solo 前，先写下主题里最想发展的两三个音；让每一次独奏都回到这个动机，听众才跟得上。

\subsection{《Change of the Guard》：把爵士带进当代编制}
\label{sec:jazz-case-change-of-guard}

Kamasi Washington 的《Change of the Guard》收录于 2015 年专辑《The Epic》。它把硬波普的语言、放克的律动与当代的制作放在同一个编制里，是年轻听众进入爵士的一扇门。

\begin{itemize}
\item \textbf{编制}：萨克斯、小号、长号前台，电声键盘/贝斯与鼓构成的节奏组，兼有嘻哈/放克的节奏感。
\item \textbf{层次}：把大乐队的声部思维与当代的叠轨结合，厚而不乱。
\item \textbf{律动}：既摇摆又带放克重音，说明爵士律动可以与当代节奏共存。
\item \textbf{编配逻辑}：保留爵士的对话与独奏内核，外层换成年轻人熟悉的音色与能量。
\end{itemize}

可借鉴：爵士不是“只能老派”。把对话与即兴留下，换一套你熟悉的音色与律动，反而更容易写出自己的声音。

\section{夜半骨架：八小节爵士三重奏}
\label{sec:jazz-yeban}

下面是一套能在夜半里直接重建的八小节骨架，用“钢琴三重奏 + 一件前台旋律”把爵士对话跑一遍。

\begin{longtable}{p{0.10\textwidth}p{0.22\textwidth}p{0.58\textwidth}}
\toprule
小节 & 段落 & 编配动作 \\
\midrule
\endfirsthead
\toprule
小节 & 段落 & 编配动作 \\
\midrule
\endhead
1—4 & Head（主题） & 钢琴：右手写 8 小节主题动机，左手只在 1/3 拍给和弦；贝斯：每拍一音的 Walking Line（根音—经过音—五度—经过音）；鼓：Ride 八分摇摆、军鼓轻点 2/4；小号/萨克斯：与钢琴主题做问答，句尾留两拍 \\
5—8 & Solo（独奏） & 前台旋律轨在 ii–V–I 上发展第 1—4 小节的动机；钢琴 Comping 降密度，只在关键拍给和弦；贝斯保持行走、不加花；鼓把军鼓换成 Ride 与轻踩镲，给独奏让路；第 8 小节末回到主题动机 \\
\bottomrule
\end{longtable}

操作顺序：

\begin{sloppypar}
\begin{enumerate}
\item 新建中等速度（约 120 BPM）、4/4 的摇摆工程，按上表建轨并命名：Melody（MIDI）、Piano（MIDI）、Bass（MIDI）、Drums（MIDI）、Master 总线。
\item 先写贝斯的八小节 Walking Line：每拍一音，用经过音连接和弦根音；Solo 听一遍，确认它在“走”而不只是点根音。
\item 写鼓：Ride 走摇摆八分，军鼓轻点 2/4，底鼓稀疏；开量化摆动，让八分真的有“长—短”。
\item 写主题：目标音落在 ii–V–I 的和弦音上，句尾留两拍；让主题有一个可发展的两音动机。
\item 写钢琴 Comping：只在关键拍给和弦，其余留空，像在回答前台旋律。
\item Solo 段把钢琴和弦减少、鼓降密度；用 F7 打开钢琴卷帘，把主题动机移到新的和弦上重复与发展。
\item 用 F3 打开设备堆叠，确认前台旋律与钢琴不抢同一音区；用 F8 看电平，低频只有贝斯一个主角。
\item 第 8 小节末回到主题动机，保存工程并导出一次八小节试听；能听出“问—答—发展—回归”就成立。
\end{enumerate}
\end{sloppypar}

\section{练习：三条路，不同取舍}
\label{sec:jazz-practice}

\begin{exercise}
\textbf{同一主题，三种独奏起手式。} 写一段 8 小节主题，再做三次 Solo：A 只用一个动机做节奏变形；B 把动机移到每个和弦上；C 先跑和弦音再回主题。比较“连贯、惊喜、可听懂”三项。A 最连贯但可能单调；B 最有结构感；C 最容易被听懂但容易像练习曲。
\end{exercise}

\begin{exercise}
\textbf{把 Swing 写出来，再改回去。} 先在夜半里把一段八分写成平均时值，播放；再打开量化摆动（或手动把后半拍后移），播放。记录哪一种“想跟着晃”。最后试着让贝斯与 Ride 的摇摆错开一点点，听弹性从哪里来。
\end{exercise}

\begin{exercise}
\textbf{四轨大乐队。} 只用四条 MIDI 轨模拟 Big Band 的声部感：把四条轨分成“问句组”和“答句组”，轮流演奏同一动机的不同音区。比较“全员齐奏”与“分组问答”哪一版更有画面；不要用更多轨，只用分工。
\end{exercise}

\section{章末扩展阅读}
\label{sec:jazz-reading}

\begin{sloppypar}
\begin{itemize}
\item \textbf{Wikipedia, “Jazz”}：\url{https://en.wikipedia.org/wiki/Jazz}。用流派条目建立爵士分支的坐标系，适合先看清全貌。
\item \textbf{Wikipedia, “Kind of Blue”}：\url{https://en.wikipedia.org/wiki/Kind_of_Blue}。Miles Davis 1959 年专辑条目，含《So What》的创作背景与调式手法；听官方音频时用艺人官方频道检索曲名。
\item \textbf{Wikipedia, “Take the 'A' Train”}：\url{https://en.wikipedia.org/wiki/Take_the_%22A%22_Train}。Duke Ellington/Billy Strayhorn 作品的条目，交代大乐队编制与年代。
\item \textbf{Wikipedia, “My Favorite Things (John Coltrane album)”}：\url{https://en.wikipedia.org/wiki/My_Favorite_Things_(John_Coltrane_album)}。Coltrane 1961 年专辑条目，说明标题曲怎样把百老汇歌曲改成调式长独奏。
\item \textbf{Wikipedia, “The Epic (Kamasi Washington album)”}：\url{https://en.wikipedia.org/wiki/The_Epic_(album)}。2015 年专辑条目，含《Change of the Guard》的编制与发行资料。
\end{itemize}
\end{sloppypar}

以上链接在 2026 年 11 月 27 日按条目名与发行资料逐条核对（专辑/作曲/年代见条目正文）；维基条目地址稳定，若所在地区无法访问，请以同名条目在当地语言版本检索。音色与独奏请到对应艺人的官方频道按曲名收听。

\section{本章要点}
\label{sec:jazz-takeaways}

\begin{itemize}
\item 爵士的核心是“对话”：主题先说、别人回答、节奏组托底、独奏发展。
\item 摇摆是演奏感，不是把音符写成长短；Ride 镲与时间感比音阶更重要。
\item 编制记住“节奏组 + 前台”：节奏组管时间和和声，前台管旋律与故事。
\item ii–V–I 与调式提供两种空间：前者练跟上变化，后者练发展动机。
\item 主题要短、可记、留一个可发展的动机；独奏要回到主题，而不是跑音阶。
\item 独奏进入时节奏组降密度，问答要真的留空，低频只留行走贝斯一个主角。
\item 夜半骨架按“贝斯—鼓—主题—Comping—Solo—音区/电平检查—回归”的顺序落地。
\end{itemize}

\begin{helpmapping}
\begin{longtable}{p{0.45\textwidth}p{0.45\textwidth}}
\toprule
小节 & 帮助锚点 \\
\midrule
什么是爵士 & \texttt{arrangement-genre-jazz-what} \\
从新奥尔良到全球 & \texttt{arrangement-genre-jazz-history} \\
配器：节奏组与前台 & \texttt{arrangement-genre-jazz-craft} \\
四首曲子的配器拆解 & \texttt{arrangement-genre-jazz-cases} \\
夜半骨架 & \texttt{arrangement-genre-jazz-yeban} \\
练习 & \texttt{arrangement-genre-jazz-practice} \\
章末扩展阅读 & \texttt{arrangement-genre-jazz-reading} \\
\bottomrule
\end{longtable}
\end{helpmapping}
